% Propietario conservado: /Users/ruben/Documents/New project/output/REVISION_ARGUMENTAL_APERTURAS_CIERRES_20260915/VII/delivery/MANUSCRIPT_VII_ES_EN_COMPLETE/source_es/manuscrito/32_eliminacion_torsional_no_lineal.tex
% FORMALIZACION_NUEVA / CERTIFICADO_NUEVO, 10-09-2026.
% Composición de las operaciones efectivas ya reunidas en 30e y 31.
% No modifica la selección de rutas, la frontera, los acoplamientos ni
% la normalización de la acción; no evalúa una amplitud cosmológica s_0.
% Propietarios y alcance: desarrollo/COMPOSICION_METRICA_REV06.md.
% Control independiente: pruebas/verificar_eliminacion_torsional_no_lineal.py.

\section{Eliminación torsional y respuesta métrica de la extensión mínima}
\label{vii:vii:sec:eliminacion-torsional-no-lineal}

La extensión mínima proporciona un campo, un peso de respuesta y una
corriente mediante una sola operación variacional. Su incidencia
transportada depende de la conexión a través de los productos de ruta de
\eqref{vii:vii:eq:derivada-incidencia-rutas}. Esa dependencia se conserva al
componer la corriente con la ecuación de Cartan--Holst. La eliminación
de la conexión se realiza, por tanto, sobre la acción completa de esa
extensión. El caso afín de
\ref{vii:vii:thm:respuesta-cuadratica} queda contenido como una
especialización exacta de la construcción siguiente.

\subsection{Composición variacional en un nivel celular}
\label{vii:vii:subsec:composicion-torsional-celular}

Fijemos un nivel finito, sus rutas, la representación de transporte y
el emparejamiento celular de
\ref{vii:vii:thm:componentes-minimo-conexion}. Sean \(p\) coordenadas
reales de la cotetrada y de los campos materiales en una trivialización
local declarada, y \(k\in\mathbb R^d\) las coordenadas de la
contorsión \(K\). La familia \(p\mapsto e(p)\) conserva la
no degeneración de la cotetrada. La conexión que actúa en cada ruta es
\begin{equation}
 \omega(p,k)=\mathring\omega(e(p))+K(k),\qquad
 F(p,k)=S_{\min}[e(p),\omega(p,k),\Psi(p)]
       =\mathfrak a(p,k)E(p,k).
 \label{vii:vii:eq:accion-minima-en-contorsion}
\end{equation}
Los espacios y las bases de esta expresión se transportan a espacios
fijos antes de derivar. Se trabaja en un abierto donde la incidencia
\(\mathsf C(p,k)\) es sobreyectiva y las familias son \(C^2\).
Esta regularidad adicional permite examinar el Hessiano de la acción.
Los acoplamientos \(G_{\rm int},c_{\rm int},\gamma_{\rm HMT}\)
conservan la carta y la convención de variación de
\ref{vii:vii:subsec:dominio-corriente-cartan}.

Sea \(Q(p,k)\) la integral, o la realización celular declarada, de
la forma cuadrática
\eqref{vii:vii:eq:forma-cuadratica-contorsion}. En bases reales,
\begin{equation}
 Q(p,k)=\tfrac12 k^{\mathsf T}\mathsf A(p)k,
 \qquad \mathsf A(p)^{\mathsf T}=\mathsf A(p).
 \label{vii:vii:eq:cuadratica-celular-contorsion}
\end{equation}
El operador \(\mathsf A\) procede de la acción gravitatoria y de
su emparejamiento, mientras que el Gramiano
\(\mathsf C\mathsf C^*\) procede del problema de extensión.
La positividad del segundo garantiza la existencia del campo
minimizante; la forma gravitatoria conserva su propia signatura.
Los términos de borde se mantienen según
\eqref{vii:vii:eq:borde-contorsion}; las fórmulas interiores siguientes
se aplican con variaciones admisibles que anulan su contribución, o
con la completación de borde correspondiente.

Definamos el covector de respuesta mediante
\begin{equation}
 d_kF(p,k)[h]= -\tfrac12 h^{\mathsf T}s(p,k),
 \qquad s(p,k)=\mathsf M_{\rm cel}(p,k)\sigma(p,k).
 \label{vii:vii:eq:covector-torsional-minimo}
\end{equation}
Aquí \(s\) y \(\sigma\) conservan los tipos distintos de
\eqref{vii:vii:eq:componentes-minimo-conexion}: el primero contiene los
coeficientes del diferencial y el segundo las coordenadas de la
corriente en el emparejamiento material. Si se cambia de base en
\(k\), también se transporta ese emparejamiento. Al diferenciar
\(s=\mathsf M_{\rm cel}\sigma\), se incluyen las derivadas de
ambos factores cuando el emparejamiento depende del parámetro.

\begin{proposition}[Ecuación compuesta de la extensión y la contorsión]
\label{vii:vii:prop:ecuacion-torsional-no-lineal}
La estacionariedad interior de la acción respecto de la contorsión es
\begin{equation}
 \mathcal R(p,k):=\mathsf A(p)k-\tfrac12s(p,k)=0.
 \label{vii:vii:eq:ecuacion-torsional-no-lineal}
\end{equation}
Cada componente de \(s\) se calcula a partir de la misma extensión:
\begin{equation}
 \begin{split}
 s_j={}&4\mathfrak a\chi\operatorname{Re}
       \langle y,(\partial_{k_j}\mathsf C)f_{\min}
                         -\partial_{k_j}b\rangle\\
 &-2\mathfrak a(\partial_{k_j}\chi)\langle b,y\rangle
       -2(\partial_{k_j}\mathfrak a)E.
 \end{split}
 \label{vii:vii:eq:seleccion-variacional-corriente-k}
\end{equation}
Cuando el emparejamiento celular realiza la ecuación de conexión de
\ref{vii:vii:thm:ecuacion-cartan-holst}, esta ecuación equivale a componer
\(\sigma(p,k)\) con la inversa de Holst, la inversa de
\(\mathsf L_e\) y la reconstrucción de contorsión:
\begin{equation}
 K=\mathsf C_e^{\rm tor}\,\mathsf L_e^{-1}
           \bigl(\kappa c\,\mathsf P_\gamma^{-1}
             \sigma[e,\mathring\omega(e)+K,\Psi]\bigr).
 \label{vii:vii:eq:composicion-implicita-cartan}
\end{equation}
El símbolo \(\mathsf C_e^{\rm tor}\) designa exclusivamente la
aplicación \(T\mapsto K\) de
\eqref{vii:vii:eq:contorsion-componentes}, con cotetrada fija.
\end{proposition}

\begin{proof}
El diferencial de \(Q+F\), en la dirección arbitraria \(h\), es
\[
 d_k(Q+F)[h]
   =h^{\mathsf T}\bigl(\mathsf A k-\tfrac12s\bigr).
\]
Su anulación para toda dirección prueba
\eqref{vii:vii:eq:ecuacion-torsional-no-lineal}.
El teorema~\ref{vii:vii:thm:componentes-minimo-conexion}, aplicado a
\(\partial_{k_j}\), da
\eqref{vii:vii:eq:seleccion-variacional-corriente-k}. Con \(p\) fijo,
\(\delta\omega=\delta K\). La ecuación de Cartan se transforma,
sucesivamente, mediante \eqref{vii:vii:eq:inversa-holst},
\eqref{vii:vii:eq:torsion-explicita} y
\eqref{vii:vii:eq:contorsion-componentes}; cada una de estas operaciones
es invertible en su dominio. Esto prueba la equivalencia con
\eqref{vii:vii:eq:composicion-implicita-cartan} cuando ambas expresiones
usan la misma realización material y celular.
\end{proof}

La composición conserva la conexión en los dos miembros. Las rutas
que recorren varias celdas pueden acoplar las componentes de la
corriente entre ellas. La reconstrucción puntual de \(T\) a partir
de una corriente dada y la resolución conjunta de
\eqref{vii:vii:eq:ecuacion-torsional-no-lineal} son, en ese caso,
operaciones sucesivas. Esta formulación finita mantiene explícita
la representación de rutas que interviene en el funcional.

\subsection{Existencia local de una rama estacionaria}
\label{vii:vii:subsec:rama-torsional-local}

\begin{theorem}[Continuación local de la solución torsional]
\label{vii:vii:thm:rama-torsional-local}
Supongamos que \((p_0,k_0)\) pertenece al abierto anterior,
\(\mathcal R(p_0,k_0)=0\), y que el Hessiano total
\begin{equation}
 \mathsf H_0
 =\mathsf A(p_0)+\partial_k^2F(p_0,k_0)
 =\mathsf A(p_0)-\tfrac12\partial_ks(p_0,k_0)
 \label{vii:vii:eq:hessiano-total-torsional}
\end{equation}
es invertible. Existen entornos de \(p_0\) y \(k_0\) en los que
una única rama \(C^1\), \(k=k_*(p)\), satisface
\eqref{vii:vii:eq:ecuacion-torsional-no-lineal} y pasa por \(k_0\).
En esa rama,
\begin{equation}
 \partial_{p_i}k_*
 =-\bigl[\mathsf A+\partial_k^2F\bigr]^{-1}
     \left[(\partial_{p_i}\mathsf A)k_*
                      +\partial_{p_i}\partial_kF\right].
 \label{vii:vii:eq:derivada-rama-torsional}
\end{equation}
\end{theorem}

\begin{proof}
La inversión del Gramiano es \(C^2\) en el abierto de rango
completo. Por ello \(F\) es \(C^2\) y \(\mathcal R\) es
\(C^1\). Su derivada respecto de \(k\), en \((p_0,k_0)\), es
\(\mathsf H_0\). El teorema de la función implícita proporciona
la rama, su regularidad y su unicidad en un entorno del punto.
También puede verse la existencia mediante la aplicación
\[
 \mathcal T_p(k)=k-\mathsf H_0^{-1}\mathcal R(p,k).
\]
Su derivada respecto de \(k\) se anula en \((p_0,k_0)\).
En una bola suficientemente pequeña tiene norma a lo sumo
\(q<1\). Reduciendo el entorno de \(p_0\), se consigue
\(\|\mathcal T_p(k_0)-k_0\|<(1-q)r\), donde \(r\) es el
radio de la bola. Así, \(\mathcal T_p\) lleva la bola cerrada
en sí misma y es contractiva. Su punto fijo es la solución local.
Diferenciar \(\mathcal R(p,k_*(p))=0\) y despejar
\(\partial_{p_i}k_*\) da
\eqref{vii:vii:eq:derivada-rama-torsional}.
\end{proof}

La hipótesis se verifica sobre la acción compuesta: el rango del
Gramiano asegura el mínimo de pantalla y el Hessiano total controla
la rama torsional. La unicidad afirmada es local alrededor de la
solución especificada; otras ramas pueden coexistir en el mismo
dominio de rango completo.

\subsection{Acción efectiva y contribución no afín}
\label{vii:vii:subsec:accion-efectiva-no-afin}

\begin{theorem}[Identidad exacta de eliminación]
\label{vii:vii:thm:eliminacion-torsional-no-lineal}
Sea \(k_*(p)\) una rama estacionaria. Supongamos, además, que el
segmento \(t\mapsto (p,tk_*)\), \(0\leq t\leq1\), permanece
en el abierto donde está definido \(F\). La corrección de acción
respecto de la realización de Levi--Civita es
\begin{equation}
 \begin{split}
 \Delta S(p)
  &:=Q(p,k_*)+F(p,k_*)-F(p,0)\\
  &=\tfrac14 k_*^{\mathsf T}s(p,k_*)
     -\tfrac12\int_0^1
           k_*^{\mathsf T}s(p,tk_*)\,dt.
 \end{split}
 \label{vii:vii:eq:accion-efectiva-no-lineal}
\end{equation}
Equivalentemente,
\begin{equation}
 \Delta S=-\tfrac14 k_*^{\mathsf T}s(p,k_*)
  +\tfrac12\int_0^1 k_*^{\mathsf T}
       \bigl[s(p,k_*)-s(p,tk_*)\bigr]\,dt.
 \label{vii:vii:eq:correccion-no-afin-integrada}
\end{equation}
Estas identidades se refieren a la acción integrada o celular.
En una realización con densidad local, se aplican a la densidad
con su emparejamiento de formas correspondiente. La acción efectiva
conserva separadamente el término de borde declarado.
\end{theorem}

\begin{proof}
Por la regla de la cadena y
\eqref{vii:vii:eq:covector-torsional-minimo},
\[
 \frac{d}{dt}F(p,tk_*)
      =-\tfrac12 k_*^{\mathsf T}s(p,tk_*).
\]
Integrar entre \(0\) y \(1\) da la diferencia de acciones
materiales. La ecuación de estacionariedad, evaluada en la
dirección \(k_*\), y la homogeneidad cuadrática de \(Q\) dan
\[
 2Q(p,k_*)=\tfrac12 k_*^{\mathsf T}s(p,k_*).
\]
La suma de ambas identidades demuestra
\eqref{vii:vii:eq:accion-efectiva-no-lineal}; sumar y restar
\(\tfrac12k_*^{\mathsf T}s(p,k_*)\) proporciona
\eqref{vii:vii:eq:correccion-no-afin-integrada}.
\end{proof}

\begin{corollary}[Especialización afín]
\label{vii:vii:cor:recuperacion-eliminacion-afin}
Si \(s(p,k)=s(p)\), entonces
\[
 \Delta S=-\tfrac14 k_*^{\mathsf T}s(p)=-Q(p,k_*).
\]
Es la realización integrada de
\eqref{vii:vii:eq:accion-efectiva-spin}.
\end{corollary}
\begin{proof}
La diferencia dentro de la integral de
\eqref{vii:vii:eq:correccion-no-afin-integrada} se anula. La segunda
igualdad procede de la identidad radial de estacionariedad.
\end{proof}

\subsection{Diferencial métrico de la acción eliminada}
\label{vii:vii:subsec:diferencial-metrico-no-lineal}

\begin{theorem}[Respuesta de las variables geométricas]
\label{vii:vii:thm:envolvente-metrica}
Sobre la rama del teorema~\ref{vii:vii:thm:rama-torsional-local},
\begin{equation}
 \partial_{p_i}\Delta S
 =\tfrac12 k_*^{\mathsf T}(\partial_{p_i}\mathsf A)k_*
  +(\partial_{p_i}F)(p,k_*)-(\partial_{p_i}F)(p,0),
 \label{vii:vii:eq:envolvente-metrica-no-lineal}
\end{equation}
donde las derivadas parciales mantienen fijas las coordenadas \(k\).
En productos hermíticos fijos, cada derivada material se calcula por
\begin{equation}
 \begin{split}
 \partial_{p_i}F={}&(\partial_{p_i}\mathfrak a)E
       +\mathfrak a(\partial_{p_i}\chi)\langle b,y\rangle\\
   &+2\mathfrak a\chi\operatorname{Re}
       \langle y,\partial_{p_i}b
                    -(\partial_{p_i}\mathsf C)f_{\min}\rangle.
 \end{split}
 \label{vii:vii:eq:diferencial-geometrico-minimo}
\end{equation}
Si \(p_i\) varía la cotetrada, \(\partial_{p_i}\mathsf C\)
incluye la variación de
\(\mathring\omega(e(p))+K(k)\), de las rutas realizadas y de
las trivializaciones utilizadas, según sus dependencias declaradas.
\end{theorem}

\begin{proof}
La derivada total de \(Q(p,k_*)+F(p,k_*)\) contiene, además
de sus derivadas parciales, el término
\[
 \bigl[\partial_kQ(p,k_*)+\partial_kF(p,k_*)\bigr]
                      \partial_{p_i}k_*.
\]
El corchete es cero por la ecuación de conexión. Restar la derivada
de \(F(p,0)\) prueba
\eqref{vii:vii:eq:envolvente-metrica-no-lineal}.
La regla del producto y
\eqref{vii:vii:eq:diferencial-minimo-completo} prueban
\eqref{vii:vii:eq:diferencial-geometrico-minimo}. La composición
\eqref{vii:vii:eq:accion-minima-en-contorsion} obliga a diferenciar
la conexión de Levi--Civita al variar \(e\) con \(k\) fijo.
\end{proof}

Cuando el producto positivo de extensión también varía, su
contribución se conserva de forma explícita. En coordenadas, sea
\(\mathsf h(p,k)>0\) su matriz hermítica y considérese el mínimo
de \(f^*\mathsf h f\) bajo \(\mathsf C f=b\). Entonces
\begin{equation}
 \mathsf L=\mathsf C\mathsf h^{-1}\mathsf C^*,\quad
 y=\mathsf L^{-1}b,\quad
 f_{\min}=\mathsf h^{-1}\mathsf C^*y,\quad
 E=\chi b^*y.
 \label{vii:vii:eq:minimo-producto-variable}
\end{equation}
Aquí \(y\) representa el covector de la restricción en bases
fijas. La matriz positiva \(\mathsf h\) conserva su papel
energético, distinto del de la métrica lorentziana.

\begin{lemma}[Variación del producto de extensión]
\label{vii:vii:lem:producto-variable-extension}
Con las convenciones anteriores,
\begin{equation}
 \delta E=(\delta\chi)b^*y
   +2\chi\operatorname{Re}\bigl[y^*(\delta b-\delta\mathsf C f_{\min})\bigr]
   +\chi f_{\min}^*(\delta\mathsf h)f_{\min}.
 \label{vii:vii:eq:diferencial-producto-variable}
\end{equation}
\end{lemma}
\begin{proof}
La fórmula \(\delta\mathsf h^{-1}
=-\mathsf h^{-1}(\delta\mathsf h)\mathsf h^{-1}\) da
\[
 \delta\mathsf L=(\delta\mathsf C)\mathsf h^{-1}\mathsf C^*
 +\mathsf C\mathsf h^{-1}(\delta\mathsf C)^*
 -\mathsf C\mathsf h^{-1}(\delta\mathsf h)
                              \mathsf h^{-1}\mathsf C^*.
\]
Al diferenciar \(E=\chi b^*\mathsf L^{-1}b\), el término
\(-\chi y^*(\delta\mathsf L)y\) produce los dos términos
con \(\delta\mathsf C\) y el término positivo con
\(\delta\mathsf h\). Esto prueba la identidad completa.
\end{proof}

En una realización material covariante cuya variación admita el
emparejamiento métrico de
\eqref{vii:vii:eq:fuentes-normalizadas-accion}, la corrección determina
\begin{equation}
 \delta_g\Delta S
 =-\tfrac12\int_M\operatorname{vol}_e\,
          \mathcal U^{S,\min}_{\mu\nu}\,\delta g^{\mu\nu},
 \qquad U^{\rm phys,\min}_{\mu\nu}
                      =c\mathcal U^{S,\min}_{\mu\nu}.
 \label{vii:vii:eq:fuente-metrica-minimo-no-lineal}
\end{equation}
Las ecuaciones \eqref{vii:vii:eq:envolvente-metrica-no-lineal} y
\eqref{vii:vii:eq:diferencial-geometrico-minimo}, completadas por
\eqref{vii:vii:eq:diferencial-producto-variable} cuando corresponde,
evalúan sus coeficientes en las variaciones métricas declaradas.
El mínimo de pantalla, la solución torsional y la derivada métrica
son tres operaciones consecutivas sobre la misma acción.
La acción reducida conserva los posibles acoplamientos entre
celdas presentes en las rutas originales.

La especialización homogénea de
\ref{vii:vii:sec:dilucion-rebote} utiliza además la realización del
tensor obtenido: su simetría espacial, el signo de la contracción
con el observador y la ley de escala. Cuando esa evaluación da
\(\rho_{\rm corr}(a)=-s_0a^{-6}\), la amplitud se lee en el
estado de referencia como
\(s_0=-a_{\rm ref}^{6}\rho_{\rm corr}(a_{\rm ref})\).
Esta lectura conserva el estado material y su normalización; el
teorema de eliminación proporciona la acción que debe contraerse.

\subsection{Modelo exacto de comprobación no afín}
\label{vii:vii:subsec:control-no-afin}

El siguiente modelo finito comprueba los coeficientes y la
dependencia no afín de las identidades. Sus parámetros pertenecen
al ejemplo algebraico y mantienen separado el problema de
realización material de las rutas.

\begin{proposition}[Eliminación exacta de una extensión de rango uno]
\label{vii:vii:prop:ejemplo-no-afin}
Sean \(H=\mathbb R^2\), \(Q=\mathbb R\),
\(\mathsf C(k)=(1\ \ k)\), \(b=1\), \(\chi=1\),
\(\mathfrak a=\lambda>0\), y \(Q_{\rm ex}(k)=k^2/2\).
Entonces
\begin{equation}
 f_{\min}=\frac{(1,k)^{\mathsf T}}{1+k^2},\quad
 F(\lambda,k)=\frac{\lambda}{1+k^2},\quad
 s(\lambda,k)=\frac{4\lambda k}{(1+k^2)^2},\quad
 \mathcal R(\lambda,k)
    =k\left(1-\frac{2\lambda}{(1+k^2)^2}\right).
 \label{vii:vii:eq:ejemplo-minimo-no-afin}
\end{equation}
Para \(\lambda=25/2\), las soluciones estacionarias son
\(0,2,-2\). En \(k_*=2\),
\begin{equation}
 \mathsf H_*=\frac{16}{5},\quad s_*=4,\quad
 Q_{\rm ex}(k_*)=2,\quad F(\lambda,k_*)=\frac52,
 \quad \Delta S=-8.
 \label{vii:vii:eq:valores-ejemplo-no-afin}
\end{equation}
La expresión afín \(-k_*s_*/4=-2\) difiere de la corrección
exacta. En la rama positiva,
\begin{equation}
 \partial_\lambda\Delta S=-\frac45,
 \qquad \partial_\lambda k_* =\frac1{20}
 \quad\text{en }\lambda=25/2.
 \label{vii:vii:eq:derivadas-ejemplo-no-afin}
\end{equation}
\end{proposition}

\begin{proof}
El Gramiano es \(1+k^2>0\), por lo que el mínimo tiene las
expresiones indicadas. Diferenciar \(F\) da
\(s=-2\partial_kF\) y después \(\mathcal R=k-s/2\).
Para \(\lambda=25/2\), la ecuación es
\(k[(1+k^2)^2-25]=0\); como \(1+k^2>0\), sus raíces
son precisamente \(0,\pm2\). El Hessiano total es
\[
 1-\frac{2\lambda(1-3k^2)}{(1+k^2)^3},
\]
que vale \(16/5\) en ambas raíces no nulas. La corrección es
\(2+5/2-25/2=-8\). En la identidad radial,
\[
 \int_0^1 k_*s(\lambda,tk_*)\,dt
    =\int_0^1\frac{200t}{(1+4t^2)^2}\,dt=20,
\]
pues una primitiva es \(-25/(1+4t^2)\).
Así, \(\Delta S=8/4-20/2=-8\), incluido el término no afín.
Para \(\lambda>1/2\), la rama positiva satisface
\(k_*^2=\sqrt{2\lambda}-1\) y
\(\Delta S=\sqrt{2\lambda}-1/2-\lambda\).
Sus derivadas dan \eqref{vii:vii:eq:derivadas-ejemplo-no-afin}.
La fórmula de envolvente proporciona el mismo resultado:
\((\partial_\lambda F)(\lambda,k_*)
 -(\partial_\lambda F)(\lambda,0)=1/5-1=-4/5\).
\end{proof}
